\begin{trapbox}
\par\noindent\begin{minipage}{\linewidth}
\begin{center}
\begin{tikzpicture}[x=1.05cm,y=1.05cm,>=stealth,font=\small]
\path[use as bounding box] (-2.5,-2.7) rectangle (7.5,3.8);
\begin{scope}
\node at (0,3.3) {对称切点：$MN$ 与 $x$ 轴重合};
\draw[->,gray] (-2,0)--(2.2,0) node[right] {$x$};
\draw[->,gray] (0,-2.4)--(0,2.5) node[above] {$y$};
\draw[blue,thick,domain=-2.3:2.3,samples=70] plot ({\x*\x/4},\x);
\draw[orange,thick] (-1.4,-.4)--(1.3,2.3) (-1.4,.4)--(1.3,-2.3);
\draw[gray,thick] (1,-2)--(1,2);
\draw[red,very thick] (-1,0)--(1,0);
\fill (1,2) circle (2pt) node[right] {$A$};
\fill (1,-2) circle (2pt) node[right] {$B$};
\fill[red] (-1,0) circle (2pt) node[below left] {$M$};
\fill[red] (1,0) circle (2pt) node[above right] {$N$};
\end{scope}
\begin{scope}[xshift=5.7cm]
\node at (0,3.3) {顶点切线：$MN\parallel x$ 轴};
\draw[->,gray] (-1.5,0)--(1.8,0) node[right] {$x$};
\draw[->,gray] (0,-2.4)--(0,2.5) node[above] {$y$};
\draw[blue,thick,domain=-2.3:2.3,samples=70] plot ({\x*\x/4},\x);
\draw[orange,thick] (0,-2.2)--(0,2.3) (-1.3,-.3)--(1.3,2.3);
\draw[gray,thick] (0,0)--(1,2);
\draw[red,very thick] (0,1)--(.5,1);
\fill (0,0) circle (2pt) node[below left] {$A=O$};
\fill (1,2) circle (2pt) node[right] {$B$};
\fill[red] (0,1) circle (2pt) node[left] {$M$};
\fill[red] (.5,1) circle (2pt) node[right] {$N$};
\node[orange,left] at (0,-1.5) {$x=0$};
\end{scope}
\end{tikzpicture}


先看两种位置：左图 $MN$ 在轴上，右图 $MN$ 离开轴；不能把重合写成两条不同直线平行。
\end{center}
\end{minipage}\par
\textbf{\textcolor{CTrap}{易错点一｜“同一个纵坐标”不总是严格意义上的“与 $x$ 轴平行”}}
\begin{itemize}[leftmargin=*,itemsep=0.25em]
\item \textbf{错误：}见到 $y_M=y_N$，不分情况就说“$MN$ 与 $x$ 轴平行”。
\item \textbf{正确：}先确认 $M\ne N$，再看共同纵坐标：不为 $0$ 时与 $x$ 轴平行；为 $0$ 时直线 $MN$ 与 $x$ 轴重合。
\item \textbf{最小反例：}$y^2=4x$ 上 $A(1,2),B(1,-2)$，得 $M(-1,0),N(1,0)$；$MN$ 就是 $x$ 轴。
\end{itemize}

\medskip
\textbf{\textcolor{CTrap}{易错点二｜忘记排除 $M=N$}}
\begin{itemize}[leftmargin=*,itemsep=0.25em]
\item \textbf{错误：}只说明纵坐标相同，就认定两点确定一条水平直线。
\item \textbf{正确：}补上 $x_N-x_M=(y_1-y_2)^2/(4p)\ne0$；其中 $A\ne B$ 保证 $y_1\ne y_2$。
\item \textbf{几何意义：}切线交点与中点高度相同，但横向一定隔着一段距离。
\end{itemize}

\medskip
\pagebreak
\textbf{\textcolor{CTrap}{易错点三｜把 $y^2=2px$ 和 $x^2=2py$ 混淆}}
\begin{itemize}[leftmargin=*,itemsep=0.25em]
\item \textbf{错误：}对 $x^2=2py$ 仍写 $y_M=y_N$。
\item \textbf{正确：}这时应写 $x_M=x_N=(x_1+x_2)/2$；$MN$ 竖直，且可能恰与 $y$ 轴重合。
\end{itemize}

\medskip
\textbf{\textcolor{CTrap}{易错点四｜误将 $2p$ 当成 $p$}}
\begin{itemize}[leftmargin=*,itemsep=0.25em]
\item \textbf{错误：}看到 $y^2=4x$ 就设 $p=4$。
\item \textbf{正确：}$2p=4$，所以 $p=2$。求 $x_M$、$|MN|$ 时尤其要检查此点。
\end{itemize}

\medskip
\textbf{\textcolor{CTrap}{易错点五｜在顶点处除以零}}
\begin{itemize}[leftmargin=*,itemsep=0.25em]
\item \textbf{错误：}直接把切线斜率写成 $p/y_1$，遇到 $y_1=0$ 就没有办法。
\item \textbf{正确：}始终使用 $yy_1=p(x+x_1)$，顶点 $(0,0)$ 的切线就是 $x=0$；本题不必求导或求切线斜率。
\end{itemize}

\medskip
\textbf{\textcolor{CTrap}{易错点六｜误以为所有弦都能用普通斜率}}
\begin{itemize}[leftmargin=*,itemsep=0.25em]
\item \textbf{错误：}对 $A(1,2),B(1,-2)$ 计算 $(y_2-y_1)/(x_2-x_1)$，分母为 $0$。
\item \textbf{正确：}用 $x=ty+q$ 表示弦，适用于所有不同点组成的弦，包括竖直弦；此时 $t=0$。
\end{itemize}

\medskip
\textbf{\textcolor{CTrap}{易错点七｜把 $p<0$ 漏掉}}
\begin{itemize}[leftmargin=*,itemsep=0.25em]
\item \textbf{错误：}默认 $p$ 必须为正，或者以为 $|MN|=(y_1-y_2)^2/(4p)$ 总为正。
\item \textbf{正确：}结论只要求 $p\ne0$；距离必须写成 $|MN|=(y_1-y_2)^2/(4|p|)$。
\end{itemize}

\medskip
\textbf{\textcolor{CTrap}{易错点八｜把结果直接套到一般椭圆、双曲线}}
该结论来自 $y^2=2px$（或交换坐标后的形式）的具体方程，不能仅凭“都是圆锥曲线”就套给椭圆和双曲线。
\end{trapbox}
